\section{Discussion}
\label{sec:discussion}

\subsection{What the universal Concentration-Gap conjecture would require}

The Concentration-Gap conjecture of
\citep[Microfoundation]{lasser2026micro} is universal:
optimization pressure on the proxy correlates with proxy-truth gap
exploitation, in any deployment context.
Theorem~\ref{thm:concgap-scoped} discharges only a \emph{scoped}
version under (REP), (DISP), (COAL), embedding compatibility
(Assumption~\ref{ass:embedding-companion}), and the
Lipschitz-transfer machinery of \citep{lasser2026micro}. The
universal conjecture would additionally require:

\begin{itemize}
\item A formal optimizer / selection model that captures
  ``optimization pressure'' as a structural property of agents
  rather than as the empirical correlate (HHI) used here. Existing
  candidates include the mesa-optimization framework of
  \citep{hubinger2019risks} and the formal-Goodhart machinery of
  \citep{elmhamdi2024goodhart, majka2025strong}, none of which
  alone yields a full proof.
\item A characterization of ``gap exploitation'' robust to the
  specific functional forms of $\Pproxy$ and $\Ttruth$.
  Theorem~\ref{thm:concgap-scoped} works at the proxy-truth-norm
  level; the universal conjecture would need to characterize the
  alignment-property-level slack $|g(\Ttruth) - g(\Pproxy)|$,
  which requires inverse-Lipschitz structure on $g$ not currently
  available in the GFM apparatus.
\end{itemize}

This is a research-program target rather than a technical
extension. The scoped version is achievable; a universal
model-free version is not, given current machinery.

\subsection{What the $g$-level reverse direction would require}

Theorem~\ref{thm:concgap-scoped}'s reverse direction operates at
the $\|\Pproxy - \Ttruth\|$ level: under high HHI plus separation
plus directional cancellation, the proxy-truth norm is bounded
below. The Lipschitz transfer of \citep[Theorem
2]{lasser2026micro} is forward-only, so this norm-level lower
bound does not transfer to a lower bound on $|g(\Ttruth) -
g(\Pproxy)|$.

Strengthening the reverse to the $g$-level would require an
additional non-degeneracy condition on $g$: a constant $g$
provides an immediate counterexample
($|g(T) - g(P)| = 0$ regardless of $\|P - T\|$). The natural
condition is bi-Lipschitz invertibility of $g$ on the relevant
subspace, but this is a strong restriction on the alignment
property's functional form.

For deployment-claim purposes, the norm-level reverse is
sufficient: it tells operators the regime in which non-trivial
proxy-truth distortion is structurally guaranteed, which the
detection layer (Layer 2) of \citep[Deployment Safety]{lasser2026paper10} can act
on. The $g$-level reverse would strengthen the result but is not
required for the deployment claim's operational utility.

\subsection{What the latent-coalition residual would require}

The (COAL) condition partitions counterparties whose correlations
are detected by audit; the latent-coalition residual
$\eta_\mathrm{latent}$ absorbs undetected-coordination risk as a
\emph{norm bound} on the residual error term $e_\mathrm{latent}
\in \mathcal{V}$ from
Assumption~\ref{ass:embedding-companion}: $\|e_\mathrm{latent}\|
\leq \eta_\mathrm{latent}$.

Bounding $\eta_\mathrm{latent}$ in a specific deployment requires
calibrating the audit's detection power against a worst-case
deviation magnitude: $\eta_\mathrm{latent} \leq p_\mathrm{miss}
\cdot \|\Delta\|_{\max}$ where $p_\mathrm{miss}$ is the probability
that a coalition of given size escapes detection thresholds and
$\|\Delta\|_{\max}$ is the worst-case per-counterparty deviation
in the deployment's threat model. The audit specifies the
detection thresholds, the detection-power model, and the
one-sided upper-confidence procedure in (COAL)'s procedure
(\S\ref{sec:audits}); $p_\mathrm{miss}$ and $\|\Delta\|_{\max}$
are calibrated through that procedure rather than read off raw
operating statistics, and the certification semantics this paper
inherits do not derive their values from first principles.

\subsection{Connection to existing Goodhart literature}

The Concentration-Gap selection theorem fits into the broader
formal-Goodhart literature in a specific position:
\begin{itemize}
\item It extends \citep{manheim2019categorizing}'s Adversarial Goodhart
  variant by giving an explicit upper bound on the proxy-truth
  divergence in terms of selection concentration, where Manheim
  \& Garrabrant give a qualitative taxonomy.
\item It complements \citep{elmhamdi2024goodhart}'s tail-based
  formal Goodhart by addressing the \emph{selection-induced} gap
  (where this paper sits) rather than the \emph{tail-event}
  gap (where El-Mhamdi \& Hoang sit). The two are independent
  failure modes; both should hold for the deployment claim to
  bind.
\item It fits \citep{majka2025strong}'s Strong/Weak/Benign
  Goodhart classification at the ``Weak Goodhart'' boundary:
  the proxy-truth gap is bounded under structural conditions,
  but the conditions can fail (Strong Goodhart) or be irrelevant
  (Benign Goodhart).
\end{itemize}

The selection theorem does not subsume any of these works; it
combines their perspectives into a single deployment-relevant
result.

\subsection{Connection to bounded co-evolution failure modes}

Theorem~\ref{thm:c7-corollary} replaces the
free-floating-constant assumption of bounded co-evolution with a
derivation from primitive parameters. The failure modes of
bounded co-evolution now factor through specific structural
conditions:
\begin{itemize}
\item (C5.HOEFF) failure: per-step LLR is unbounded, so
  $\Bclip$ is undefined. Detection: SPRT-applicability audit
  (paper 10 Audit 7) catches this.
\item (C7.RATE) failure: exposure event count grows faster than
  $\lammax \cdot \taumeta$. Detection: (C7.RATE) audit (this
  paper, \S\ref{sec:audits}) catches this.
\item Channel-projection failure: ledger event-classification
  policy doesn't respect the zero-outside-engagement convention.
  Detection: structural-projection audit (this paper,
  \S\ref{sec:audits}) catches this.
\end{itemize}

This decomposition is the structural payoff of
Theorem~\ref{thm:c7-corollary}: bounded-co-evolution failure is
no longer an opaque single-mode risk but a triage of three
specific structural-condition failures, each with its own audit
detection.

